\documentclass[french,11pt]{book}
\usepackage{teach}
\usepackage[french]{babel}
\usepackage{amsmath,amssymb}
\usepackage{array}
\newcommand{\R}{\mathbb{R}}
\begin{document}
\begin{center}
{\Large \textbf{Exercices -- Fonction inverse}}\\[2mm]
Terminale technologique
\end{center}

\section*{Exercices}
\begin{enumerate}
  \item Calculer les images de $-4$, $-1$, $0,5$, $2$ et $10$ par la fonction inverse.
  \item Résoudre dans $\mathbb{R}^*$ les équations $\frac{1}{x}=4$, $\frac{1}{x}=-0,2$ et $\frac{1}{x}=\frac{5}{3}$.
  \item Comparer, sans calculatrice, $\frac{1}{2}$ et $\frac{1}{5}$ puis $\frac{1}{-2}$ et $\frac{1}{-5}$. Expliquer avec les variations de la fonction inverse.
  \item Une grandeur $y$ est inversement proportionnelle à $x$ et $y=12$ quand $x=5$. Déterminer $y$ lorsque $x=8$.
\end{enumerate}

\end{document}
